\chapter{Introduction}
The human brain is one of the most complex things that exists in nature. It controls everything a person does, from walking and talking to thinking and feeling emotions. Even when a person is sitting quietly and doing nothing, the brain is still working. It never fully stops. With the growth of technology, it has become possible to study the brain in ways that were not possible before. This project is one such attempt to use modern technology to understand two important things about the brain, which are the mental state of a person and which hand a person is imagining moving.

In daily life, every person goes through different mental conditions. There are times when a person feels very fresh and focused. There are also times when a person feels mentally tired, dull, or distracted. These changes happen because of many reasons like lack of sleep, long hours of studying or working, stress, or even boredom. These different mental conditions are referred to as cognitive states. The word cognitive simply means anything related to thinking, attention, memory, or mental processing. So cognitive assessment means measuring or evaluating how the brain is functioning at a given moment. Cognitive assessment is not only useful in medical cases. It is also useful in everyday situations. For example, if a student is studying for a long time, it is helpful to know if the student is still mentally active or has become too tired to absorb new information. Similarly, in workplaces where tasks require high concentration, like air traffic control or surgery, knowing the cognitive state of the person doing the task can help avoid mistakes.

The brain produces electrical signals continuously. These signals change depending on what the brain is doing at a given moment. When a person is relaxed, the signals look different compared to when the person is focused or stressed. These signals can be measured from outside the head without any surgery or internal device. This is why brain signal based cognitive assessment has gained a lot of interest in recent years. The method used to record brain electrical signals from the scalp is called Electroencephalography, which is commonly known as EEG. In this method, small sensors called electrodes are placed on different positions on the scalp. These electrodes pick up the tiny electrical signals produced by the brain.
Along with cognitive assessment, this project looks at motor imagery. When a person only imagines moving the left or the right hand, without moving at all, the part of the brain that controls that hand still changes its activity, and this change can be picked up from electrodes placed over the motor cortex. If a computer can tell which hand was imagined, the person can control something, such as a cursor or a wheelchair, by thought alone. A longer-term aim of the work is imagined speech, where a person says a word only in their mind, and the motor imagery work is the first step towards it. 



\section{Motivation}
The brain produces electrical signals all the time but reading those signals and making useful sense out of them is not easy. EEG signals are noisy, complex, and they change a lot from person to person. Getting reliable information from them, especially when using a simple wearable headset in everyday conditions, is still a difficult problem. This problem becomes even more important to solve when looking at the current global situation. The World Health Organization estimates that over 3 billion people, more than one in three, live with a neurological condition, and that such conditions now cause more ill health and disability than any other group of diseases \cite{who2024neuro}. These numbers show that there is a real need for simple and non-invasive tools that can monitor brain activity and help people in practical ways.

Now when it comes to recording brain signals, there are two ways of doing it. The invasive way involves placing electrodes inside the skull through surgery, which gives clear signals but comes with serious risks like infection, bleeding, or permanent brain damage. Because of this, non-invasive EEG is used in this project, where electrodes sit on the outside of the scalp without any surgery. Within non-invasive EEG, gel based electrodes give better signal quality but require a trained person to apply gel on each electrode point, which is messy, slow, and not practical for everyday use. Dry electrodes on the other hand need no gel, are quick to set up, and the person can wear the headset themselves without any help. Even though they may pick up slightly more noise, they are a much better choice for a real world wearable setup like this one.

One of the things this project works on is helping people who cannot speak. There are people whose brain works completely fine but due to conditions like paralysis or nerve damage, they have no way to express what they want to say. Saying something as basic as being in pain or needing water becomes impossible for them. The brain still produces clear electrical signals when such a person only imagines a movement or a word. In this project these signals are recorded with an EEG headset, and a machine learning model is trained to tell which hand the person was imagining moving. This is the simplest form of a brain-computer interface and the first step towards decoding imagined words.

The second thing this project works on is understanding the mental condition of a person at a given moment. In daily life, whether it is a student studying for long hours, a working professional under deadline pressure, or anyone going through a mentally tiring day, people usually do not realize when their brain has stopped working well. They keep going even after their focus has already dropped. By analyzing EEG signals and calculating a cognitive score, this project offers a simple way to know the mental state of a person



\section{Objective of the Project}

The main objective of this project is to develop an EEG-based cognitive assessment system that can record and analyse brain signals for different applications. The project mainly focuses on telling left-hand from right-hand motor imagery using EEG data, and on analysing cognitive state using frequency band analysis.\\

\begin{itemize}[topsep=2pt, itemsep=4pt]
    \item To develop an EEG database for motor imagery (MI) actuation.

    \item To design a motor imagery classifier system.

    \item To compute a cognitive score using EEG frequency band analysis.
\end{itemize}
\section{Organization of the Report}
This report is divided into seven chapters that walk through the development of the EEG-based cognitive assessment model. The first chapter explains what EEG is, why it is useful in brain-computer interfaces, and what motivated this project given the growing number of people affected by neurological disorders. The second chapter looks at existing research on wearable EEG systems, public motor imagery datasets, including the stroke patient dataset and the BCI Competition IV 2b dataset, methods for feature extraction and classification, speech decoding and mental state assessment, and ends by identifying the gaps this project tries to address. The third chapter gives a broad overview of the entire system, from recording and cleaning EEG signals all the way to extracting features and using a machine learning model to generate an output. The fourth chapter covers the hardware, which is the GELO BCI headset, Raspberry Pi 5, aluminium heat sink casing and Waveshare LCD screen, and explains the role each component plays. The fifth chapter shifts to the software, detailing the tools and libraries used along with the step-by-step process for signal acquisition, preprocessing, feature extraction, classification and cognitive score calculation. The sixth chapter presents the results: the reproduction of the stroke dataset analysis, the ERD/ERS, power spectral density, C3/C4 power ratio and SVM results on the BCI Competition IV 2b dataset, the same analysis on the data recorded with the headset, and the cognitive assessment results. The seventh chapter summarises what was achieved and sets out the next steps, such as a better recording protocol, per-user calibration and deployment on the Raspberry Pi.
